% ===== 原书 p.32 =====
\section{等变 K 理论与齐性空间}
\subsection{G-向量丛与 $K_G(X)$}

下面考虑 K 理论的推广，把紧 Lie 群 $G$ 的作用计入。回忆 Lie 群是带有光滑流形结构的群，其乘法与取逆映射都光滑。特别地，每个有限群都是零维 Lie 群。所谓流形 $X$ 上的左 $G$-作用，是指光滑映射 $G\times X\to X$，记为 $(g,x)\mapsto g\cdot x$，满足 $(gg')\cdot x=g\cdot(g'\cdot x)$ 与 $e\cdot x=x$。类似地，右作用是一个映射 $(g,x)\mapsto x\cdot g$，满足 $x\cdot(gg')=(x\cdot g)\cdot g'$。两种情形下分别称 $G$ 从左侧或右侧作用在 $X$ 上。所谓 $G$-空间，是指流形 $X$ 连同指定的 $G$-作用。

\begin{definition}
称 $X$ 上的向量丛 $E$ 是\emph{G-向量丛}（G-vector bundle），若 $E$ 与 $X$ 都是 $G$-空间，且下面两个条件成立：
\begin{enumerate}[label=\arabic*.,leftmargin=2.2em]
\item 丛投影 $\pi:E\to X$ 与 $E$、$X$ 上的 $G$-作用交换。换言之，对所有 $v\in E$、$g\in G$ 有 $g\cdot\pi(v)=\pi(g\cdot v)$；
\item 对每个 $g\in G$ 与每个 $x\in X$，映射 $E_x\to E_{g\cdot x}$，$v\mapsto g\cdot v$，都是线性的。
\end{enumerate}
\end{definition}

\begin{definition}
给定 $G$-空间 $X$，把 Grothendieck 构造施于 $\Vect_G(X)$，即 $X$ 上复 $G$-向量丛的 $G$-同构类构成的半群，所得群记作 $K_G(X)$。对 $K_G(X)$ 的研究称为\emph{等变 K 理论}。
\end{definition}

§1 中关于普通 K 理论的结论都可直接搬到等变情形，只需把记号作相应改动。最明显的例外有两条：分裂原理只在 $G$ 交换时成立；等变 Thom 同构定理的证明要困难得多。等变 K 理论的完整论述见 [17]。

\begin{definition}
给定连续群同态 $\varphi:H\to G$，任一 $G$-空间 $X$ 都可看作 $H$-空间，作用取 $h\cdot x=\varphi(h)\cdot x$。把这一构造施于 $X$ 上的 $G$-向量丛，即得映射 $\varphi^*:K_G(X)\to K_H(X)$。
\end{definition}

注意若 $G$ 是平凡群，则 $K_G(X)$ 化归普通 K 理论中的环 $K(X)$。又注意若 $X$ 是一点，则 $X$ 上的复 $G$-向量丛就是复 $G$-模。此时 $K_G(\pt)$ 成为 $G$ 的表示环 $R(G)$，它由把 Grothendieck 构造施于 $G$ 的有限维复表示空间构成的半群得到。一般地，考虑由 $X$ 到一点的映射诱导的同态 $R(G)\to K_G(X)$，可知 $K_G(X)$ 是 $R(G)$-模。一个极端情形是

\begin{proposition}
若 $G$ 平凡地作用在 $X$ 上（即 $g\cdot x=x$ 对一切 $g$），则 $K_G(X)\cong K(X)\otimes R(G)$。
\end{proposition}

证明见 [17, §2.2]；有限群情形见 [6, §1.6]。另一个极端给出

\begin{proposition}
若 $G$ 自由地作用在 $X$ 上（即 $g\cdot x=x$ 当且仅当 $g=e$），则 $K_G(X)\cong K(X/G)$，其中 $X/G$ 是把 $x$ 与所有 $g\cdot x$ 等同所得的轨道空间（orbit space）。
\end{proposition}

\begin{proof}
设 $E$ 是 $X$ 上的 $G$-向量丛。$G$-作用此时平凡，故 $E/G$ 是 $X/G$ 上的普通向量丛。\footnote{还须验证 $E/G$ 局部平凡。}于是得到同态 $K_G(X)\to K(X/G)$。

% ===== 原书 p.33 =====
反过来，考虑由标准投影 $p:X\to X/G$ 诱导的映射。给定 $X/G$ 上的普通向量丛 $F$，拉回丛 $p^*F$ 的全空间为
\[
\{(x,v)\in X\times F\mid p(x)=\pi(v)\}.
\]
赋予 $G$-作用 $g\cdot(x,v)=(g\cdot x,v)$，可把 $p^*F$ 做成 $X$ 上的 $G$-向量丛，从而得到反向同态 $K(X/G)\to K_G(X)$。由典范向量丛同构 $E\to p^*(E/G)$ 与 $(p^*F)/G\to F$，结论立即成立。
\end{proof}

\transnote{本证明首句原文作「$G$-作用此时平凡」，与命题的自由作用前提相悖，应读作「自由」；正是自由性让 $E/G$ 成为 $X/G$ 上的普通向量丛。另外，同构 $K_G(X)\cong K(X)\otimes R(G)$ 只在作用平凡时成立，此处 $\otimes$ 是普通张量积，原文未加帽。}

\subsection{齐性空间}

从等变 K 理论的视角看，特别有趣的一类流形是齐性空间。设 $H$ 是 $G$ 的闭子群，考虑左陪集商空间 $G/H$，赋以商拓扑。所得空间是光滑流形（见 [1, p. 21]）。若 $G$ 传递地作用在流形 $X$ 上（transitive action，即对任意 $x,y\in X$，存在 $g\in G$ 使 $g\cdot x=y$），则称 $X$ 是\emph{齐性空间}。进一步，取 $H$ 为稳定子群（即固定某点的子群），则有同构 $X\cong G/H$。

齐性空间的一些简单例子如下：
\begin{itemize}[leftmargin=1.6em]
\item 任一 Lie 群：$G\cong G/\{e\}$；
\item 球面：$S^n\cong\SO(n+1)/\SO(n)$；
\item 奇维球面（odd sphere）：$S^{2n+1}\cong\U(n+1)/\U(n)$，或 $S^{2n+1}\cong\SU(n+1)/\SU(n)$（$n>0$）；
\item 实 Grassmann 流形：$G(n,k)\cong O(n+k)/O(n)\times O(k)$；
\item 实射影空间：$\R P^n\cong O(n+1)/O(n)\times O(1)$；
\item 复 Grassmann 流形：$G_{\C}(n,k)\cong\U(n+k)/\U(n)\times\U(k)$；
\item 复射影空间：$\C P^n\cong\U(n+1)/\U(n)\times\U(1)$；
\item 透镜空间（lens space）：$L(n,k)=S^{2n+1}/\Z_k\cong\U(n+1)/\U(n)\times\Z_k$。
\end{itemize}

由上述命题，注意 $H$ 自由地作用在 $G$ 上（标准作用 $h\cdot g=hg$），得到 $K(G/H)\cong K_H(G)$。不过我们的兴趣更在 $G/H$ 上的 $G$-向量丛。给定有限维 $H$-模 $M$，可构造乘积 $G\times_H M$，即 $G\times M$ 模去 $H$ 的作用：
\[
G\times_H M=G\times M\Big/(g,x)\sim h\cdot(g,x)=(gh^{-1},hx).
\]
设 $\pi:G\times_H M\to G/H$ 是投影映射 $(g,x)\mapsto gH$。它与 $G/H$ 及 $G\times_H M$ 上的标准 $G$-作用 $g'\cdot gH=g'gH$、$g'\cdot(g,x)=(g'g,x)$ 显然交换，因此 $G\times_H M$ 是 $G/H$ 上以 $M$ 为纤维的向量丛。事实上下面证明，$G/H$ 上任一 $G$-向量丛都可由这一构造得到。

\begin{proposition}
若 $H$ 是 Lie 群 $G$ 的闭子群，则 $K_G(G/H)\cong R(H)$。
\end{proposition}

\begin{proof}
设有 $G/H$ 上的 $G$-向量丛 $E$。$E$ 上的 $G$-作用限制为 $E_H$ 上的 $H$-作用，故 $E$ 在陪集 $H$ 处的纤维是有限维 $H$-模 $M$。$E$ 上的 $G$-作用 $G\times E\to E$ 限制为 $G$-同态 $G\times M\to E$，它在 $G\times M$ 的 $H$-作用下不变，故诱导 $G$-映射 $\alpha:G\times_H M\to E$。要证 $\alpha$ 是 $G$-同构。由 $G$-向量丛的定义，映射

% ===== 原书 p.34 =====
$E_H\to E_{gH}$，$x\mapsto g\cdot x$，对每个 $g\in G$ 都是同构（其逆为 $x\mapsto g^{-1}\cdot x$）。把 $\alpha$ 限制到任一陪集 $gH$ 处的纤维，得 $\alpha|_{gH}:gH\times_H M\to E_{gH}$，它恰是上述同构。于是 $\alpha$ 是整体的向量丛同构。
\end{proof}

\transnote{这一命题是 §4 的核心机制：$G$ 传递地作用在 $G/H$ 上，丛的整体行为由它在单点 $H$ 处的纤维（一个 $H$-模）决定，所以等变 K 理论把 $G/H$ 上的几何问题化归为 $H$ 的表示论。后文的 Bott 定理正是沿这条思路把指标化成表示环里的元素。}

\begin{example}
考虑球面 $S^2\cong\SO(3)/\SO(2)$，得
\[
K_{\SO(3)}(S^2)\cong R(\SO(2))\cong\Z[z,z^{-1}].
\]
特别地，这个环比 $K(S^2)\cong\Z^2$ 大得多。
\end{example}

用 $G$-作用在齐性空间 $G/H$ 上「平移」，可以看出等变向量丛的整体行为由它在单点的行为决定。用完全类似的论证，可证上述结果的微小推广。

\begin{proposition}
若 $X$ 是局部紧\footnote{这里用紧支撑 $K_G$ 理论，定义方式与普通情形类似。}的 Hausdorff $H$-空间，则 $K_G(G\times_H X)\cong K_H(X)$。
\end{proposition}

\begin{example}
齐性空间 $G/H$ 的切丛是实向量丛，在陪集 $H$ 处的纤维为 $\mathfrak{g}/\mathfrak{h}$，其中 $\mathfrak{g}$、$\mathfrak{h}$ 分别是 $G$、$H$ 的 Lie 代数。\footnote{Lie 群 $G$ 的 Lie 代数 $\mathfrak{g}\cong TG_e$ 是 $G$ 在单位元处的切空间，括号由左不变向量场的括号给出。}注意 $T(G/H)$ 是 $G$-向量丛，且 $\mathfrak{g}/\mathfrak{h}$ 带有自然的 $H$-作用（由映射 $A_h:gH\mapsto hgH$ 诱导的作用 $\mathrm{Ad}:H\to\Aut(\mathfrak{g}/\mathfrak{h})$），可知 $T(G/H)\cong G\times_H(\mathfrak{g}/\mathfrak{h})$。于是得到同构
\[
K_G\big(T(G/H)\big)\cong K_G\big(G\times_H(\mathfrak{g}/\mathfrak{h})\big)\cong K_H(\mathfrak{g}/\mathfrak{h}).
\]
\end{example}

\subsection{G-指标定理}

若 $E$ 是紧 $G$-空间 $X$ 上的 $G$-向量丛，则可把光滑截面空间 $\Gamma E$ 做成（无限维）$G$-模：

\begin{definition}
若 $s:X\to E$ 是光滑截面（不必与 $X$、$E$ 上的 $G$-作用交换），则给 $\Gamma E$ 赋予 $G$-作用 $(g\cdot s)(x)=g\cdot s(g^{-1}\cdot x)$。
\end{definition}

给定 $X$ 上的复 $G$-向量丛 $E$、$F$，自然要考察与 $\Gamma E$、$\Gamma F$ 上 $G$-作用交换的椭圆算子 $P:\Gamma E\to\Gamma F$。此时称 $P$ 是\emph{G-不变的}。注意 $TX$ 是 $G$-空间\footnote{设 $L_g:X\to X$ 是映射 $x\mapsto g\cdot x$，$dL_g:TX\to TX$ 是它的微分。取 $g\cdot(x,\xi)=dL_g(x,\xi)$，即得 $TX$ 上的标准 $G$-作用。}，$\pi^*E$、$\pi^*F$ 都是 $TX$ 上的 $G$-向量丛，故象征 $\sigma(P):\pi^*E\to\pi^*F$ 是 $G$-映射，从而给出 $K_G(TX)$ 中的一个元素。与前面一样，等变情形下的拓扑指标定义为合成
\[
\tind_G^X:K_G(TX)\xrightarrow{\ i_!\ }K_G(TM)\xrightarrow{\ j_!^{-1}\ }K_G(\pt)\cong R(G),
\]
其中 $i:X\to M$ 是把 $X$ 嵌入有限维实 $G$-模 $M$ 的嵌入，$j:\pt\to M$ 是原点的包含。

由于 $P$ 是椭圆的，$\Ker P$ 与 $\Coker P$ 都是有限维的。又因 $P$ 是 $G$-不变的，$\Ker P$、$\Coker P$ 在 $\Gamma E$、$\Gamma F$ 的 $G$-作用下不变。于是 $\Gamma E$、$\Gamma F$ 上的 $G$-作用诱导 $\Ker P$、$\Coker P$ 上的 $G$-作用，$\Ker P$、$\Coker P$ 都是有限维 $G$-模。据此可定义

% ===== 原书 p.35 =====
\begin{definition}
$G$-不变椭圆算子 $P$ 的 $G$-指标\footnote{$G$-指标与普通指标由等式 $\Ind P=\dim\Ind_G P$ 相联系。}为
\[
\Ind_G=[\Ker P]-[\Coker P]\in R(G).
\]
\end{definition}

利用 $G$-指标可把解析指标定义为同态
\[
\aind_G^X:K_G(TX)\to K_G(\pt)\cong R(G),
\]
由 $\aind_G^X=\Ind_G P$ 给出，其中 $P$ 是 $G$-不变椭圆算子且 $\sigma(P)=x$。§3 中给出的 Atiyah–Singer 指标定理证明可自然地搬到等变情形（作相应的必要改动：把 $\Z$ 换成 $R(G)$，各处补上 $G$）。于是得到 $\aind_G=\tind_G$，G-指标定理随即成立。

\begin{theorem}[G-指标定理]
若 $P:\Gamma E\to\Gamma F$ 是紧流形 $X$ 上的 $G$-不变椭圆（伪微分）算子，则
\[
\Ind_G P=\tind_G^X\sigma(P).
\]
\end{theorem}

\subsection{诱导表示}

本节回顾 Lie 群表示论的一些基本结论，用表示环的记号叙述。更完整的论述见 [1]。以下取 $G$ 为紧 Lie 群。称 $G$-模是\emph{不可约的}，若它没有 $G$-不变真子空间；称两个 $G$-模 $M$、$N$ \emph{等价}，若存在 $G$-同构 $M\to N$。取 $\{w_\alpha\}$ 为 $G$ 的一族完全的不等价不可约表示，则表示环 $R(G)$ 是由 $\{w_\alpha\}$ 生成的自由阿贝尔群。定义两个 $G$-模的\emph{缠绕数}（intertwining number）为
\[
\tau(M,N)=\dim_\C\Hom_G(M,N),
\]
可把它扩充为 $R(G)$ 上的内积 $\ip{\cdot}{\cdot}_G$。由 Schur 引理，若 $M$、$N$ 都不可约，则
\[
\ip{M}{N}_G=
\begin{cases}
1, & M\text{ 与 }N\text{ 等价},\\
0, & \text{其他},
\end{cases}
\]
于是 $\{w_\alpha\}$ 关于此内积正交归一（orthonormal）。

其次，形式地定义表示环的无限扩张（infinite extension）$\widehat R(G)$，它是所有可能无限的线性组合 $\sum_\alpha a_\alpha w_\alpha$（$a_\alpha\in\Z$）构成的集合。显然有自然包含 $R(G)\subset\widehat R(G)$。$R(G)$ 上的内积随即扩充为配对（pairing）$R(G)\otimes\widehat R(G)\to\Z$：
\[
\Big\langle x,\sum_\alpha a_\alpha w_\alpha\Big\rangle_G=\sum_\alpha a_\alpha\ip{x}{w_\alpha}_G,
\]
其中 $x\in R(G)$，$\sum_\alpha a_\alpha w_\alpha\in\widehat R(G)$。注意此和其实是有限和，因为除有限多个 $\alpha$ 外 $\ip{x}{w_\alpha}$ 为零。

\begin{definition}
给定连续群同态 $\varphi:H\to G$，定义形式诱导表示（formal induced representation）$\varphi_*:R(H)\to\widehat R(G)$ 为映射
\[
\varphi_*:x\mapsto\sum_\alpha\ip{\varphi^*w_\alpha}{x}_H w_\alpha,
\]
其中 $\varphi^*:R(G)\to R(H)$ 是等变 K 理论诱导的映射。
\end{definition}

\transnote{求和指标 $\alpha$ 取遍 $G$ 的不可约表示，而配对 $\ip{\varphi^*w_\alpha}{x}_H$ 在 $H$ 上进行：$x\in R(H)$，$\varphi^*w_\alpha$ 是 $w_\alpha$ 限制成 $H$-模所得。原文紧接着的伴随关系把 $x$ 印作 $x\in R(G)$，应为 $R(H)$。}

% ===== 原书 p.36 =====
把上述定义用扩充内积展开，可知诱导表示满足伴随关系
\[
\ip{y}{\varphi_*x}_G=\ip{\varphi^*y}{x}_H
\]
对所有 $x\in R(G)$、$y\in R(G)$ 成立。若 $H$ 是 $G$ 的闭子群，$i:H\to G$ 是包含映射，则可为诱导表示构造一个几何上更直观的替代定义。特别地，给定有限维 $H$-模 $M$，回忆 $G/H$ 上向量丛 $G\times_H M$ 的光滑截面空间 $\Gamma(G\times_H M)$ 具有无限维 $G$-模结构。把这一映射延拓到表示环，得到

\begin{definition}
给定闭子群 $H$ 的包含映射 $i:H\to G$，定义\emph{诱导表示} $i_*:R(H)\to\widehat R(G)$ 为映射
\[
M\mapsto\Gamma(G\times_H M)
\]
到 $R(H)$ 的延拓，其中 $M$ 是有限维 $H$-模。
\end{definition}

为说明这两个定义等价，需要证明下面的 Frobenius 定理：

\begin{theorem}[Frobenius 互反]
若 $W$ 是 $G$-模，$M$ 是 $H$-模，则
\[
\Hom_G(W,i_*M)\cong\Hom_H(i^*W,M),
\]
其中 $i_*M=\Gamma(G\times_H M)$，$i^*W$ 是 $W$ 限制为 $H$-模所得。
\end{theorem}

\begin{proof}
设有 $G$-不变同态 $F:W\to\Gamma(G\times_H M)$。把它与映射 $\Gamma(G\times_H M)\to M$，$s\mapsto s(H)$，复合，得到相应的 $H$-同态 $f:W\to M$。注意对 $h\in H$、$w\in W$ 有
\[
\begin{aligned}
f(h\cdot w)&=F(h\cdot w)(H)=(h\cdot F(W))(H)\\
&=h\cdot\big(F(w)(h^{-1}H)\big)=h\cdot\big(F(w)(H)\big)=h\cdot f(w).
\end{aligned}
\]
反过来，给定 $H$-不变同态 $f:W\to M$，考虑相应的映射 $F:W\to\Gamma(G\times_H M)$，它对 $gH\in G/H$、$w\in W$ 定义为
\[
F(w)(gH)=\big(g,f(g^{-1}\cdot w)\big)\in G\times_H M.
\]
由于 $f:W\to M$ 是 $H$-不变的，此映射良定义。此外注意
\[
\begin{aligned}
F(g'\cdot w)(gH)&=\big(g,f(g^{-1}g'\cdot w)\big)
=g'\cdot\Big(g'^{-1}g,f\big((g'^{-1}g)^{-1}\cdot w\big)\Big)\\
&=g'\cdot F(w)(g'^{-1}\cdot gH)=\big(g'\cdot F(w)\big)(gH)
\end{aligned}
\]
对所有 $g\in G$ 成立。故 $F:W\to\Gamma(G\times_H M)$ 是 $G$-同态。注意两个映射 $f\mapsto F$ 与 $F\mapsto f$ 互为逆，结论立即成立。
\end{proof}

\begin{corollary}
若 $i:H\to G$ 是闭子群 $H$ 在 $G$ 中的包含映射，则对任一 $H$-模 $M$ 有
\[
\big[\Gamma(G\times_H M)\big]=\sum_\alpha\ip{i^*w_\alpha}{[M]}_H w_\alpha\in\widehat R(G),
\]
于是诱导表示的两个定义一致。
\end{corollary}

\begin{proof}
由 Frobenius 互反定理，
\[
\begin{aligned}
\big[\Gamma(G\times_H M)\big]
&=\sum_\alpha\dim\Hom_G\big(w_\alpha,\Gamma(G\times_H M)\big)\cdot w_\alpha\\
&=\sum_\alpha\dim\Hom_H(i^*w_\alpha,M)\cdot w_\alpha
=\sum_\alpha\ip{i^*w_\alpha}{[M]}_H w_\alpha,
\end{aligned}
\]
这正是所要。
\end{proof}

% ===== 原书 p.37 =====
\subsection{齐性微分算子}

作为 G-指标定理的特例，下面考察齐性空间 $G/H$ 上的 $G$-不变伪微分算子\footnote{本节所有结论对伪微分算子同样成立，但为记号简便，以下省略「伪」字。} $P:\Gamma E\to\Gamma F$ 的 $G$-指标。此时称 $P$ 是\emph{齐性微分算子}。借助 Peter-Weyl 定理，此情形的分析归结为诱导映射的表示论。又因齐性算子这类对象附加的结构过多，下面很快会看到，$P$ 的指标只依赖向量丛 $E$、$F$，并不直接依赖算子本身。仿照 [12]，定义齐性象征如下：

\begin{definition}
对 $G/H$ 上的齐性算子 $P:\Gamma E\to\Gamma F$，回忆可写 $E\cong G\times_H M$、$F\cong G\times_H N$，其中 $M$、$N$ 是有限维 $H$-模。于是定义 $P$ 的\emph{齐性象征}（homogeneous symbol）为
\[
\hat\sigma(P)=[M]-[N]\in R(H).
\]
\end{definition}

现在把齐性象征与普通象征 $\sigma(P)\in K_G(T(G/H))$ 作比较。后者由 $T(G/H)$ 上的齐性复形
\[
0\longrightarrow\pi^*E\xrightarrow{\ \sigma\ }\pi^*F\longrightarrow 0
\]
表示，其中 $\sigma:\pi^*E\to\pi^*F$ 是 $G$-同态。设 $s:\pt\to T(G/H)$ 表示把恒等陪集 $H\in G/H$ 经零截面映入 $T(G/H)$ 的包含映射。诱导映射 $s^*:K_G(T(G/H))\to K_G(\pt)=R(G)$ 把 $\sigma(P)$ 映为一点上的复形\footnote{注意映射 $\sigma(H,0):M\to N$ 是同态，但一般不是 $H$-同态。}
\[
0\longrightarrow M\xrightarrow{\ \sigma(H,0)\ }N\longrightarrow 0.
\]
于是
\[
s^*\sigma(P)=[M]-[N]=\hat\sigma(P)\in R(H).
\]
现在可计算齐性微分算子的 $G$-指标。

\begin{theorem}[Bott]
若 $P$ 是 $G/H$ 上的齐性微分算子，则
\[
\Ind_G P=i_*\hat\sigma(P)=i_*\circ s^*\sigma(P)\in\widehat R(G),
\]
其中 $i_*:R(H)\to\widehat R(G)$ 是形式诱导同态。特别地，$i_*\hat\sigma(P)$ 其实是 $R(G)\subset\widehat R(G)$ 中的有限元素。
\end{theorem}

\transnote{等号两端所在的环不同，读时须小心：$s^*\sigma(P)$ 落在 $R(G)$，而 $[M]-[N]$ 是 $H$-模的形式差，落在 $R(H)$。原因在于 $s$ 把点送到 $(H,0)$，该点的迷向群是 $H$ 而非 $G$；$s$ 本身不是 $G$-映射，故 $s^*\sigma(P)$ 只带 $H$-对称性。}

由诱导表示的替代定义，
\[
i_*\big(\hat\sigma(P)\big)=i_*\big([M]-[N]\big)=\big[\Gamma(G\times_H M)\big]-\big[\Gamma(G\times_H N)\big]=[\Gamma E]-[\Gamma F],
\]
于是 Bott 定理可改写成更简洁的形式。

\begin{theorem}
若 $P:\Gamma E\to\Gamma F$ 是齐性微分算子，则
\[
\Ind_G P=[\Gamma E]-[\Gamma F]\in\widehat R(G),
\]
其中 $[\Gamma E]-[\Gamma F]$ 其实是 $R(G)\subset\widehat R(G)$ 中的有限元素。
\end{theorem}

% ===== 原书 p.38 =====
为证明此定理，须把无限维 $G$-模 $\Gamma E$ 分解成有限维 $G$-不变子空间。限制在这些子空间上，定理容易验证，再把结论提升到光滑截面的整体空间。这些子空间的自然候选如下。

\begin{definition}
对每个 $w_\alpha$，定义 $\Gamma_\alpha E$ 为同态
\[
i_\alpha:w_\alpha\otimes\Hom_G(w_\alpha,\Gamma E)\to\Gamma E
\]
的像，其中 $i_\alpha:(w\otimes\varphi)\mapsto\varphi(w)$。于是 $\Gamma_\alpha(E)$ 是 $\Gamma E$ 的有限维 $G$-不变子空间，其在表示环中的像为 $[\Gamma_\alpha E]=\ip{w_\alpha}{[\Gamma E]}_G w_\alpha\in R(G)$。
\end{definition}

注意 $\Gamma E$ 在表示环的无限扩张 $\widehat R(G)$ 中的类由 $[\Gamma E]=\sum_\alpha\ip{w_\alpha}{[\Gamma E]}_G w_\alpha$ 给出。此式提示可（至少形式地）得到无限分解 $\Gamma E=\oplus_\alpha\Gamma_\alpha E$。下述著名结果把这一分解落实。

\begin{lemma}[Peter-Weyl]
设 $L^2(E)$ 是把 $\Gamma E$ 关于某个 $G$-不变内积完备化所得的 Hilbert 空间，则诸子空间 $\Gamma_\alpha E$ 构成 $L^2(E)$ 的一族完备正交子空间。
\end{lemma}

现在已有足够准备，可较快证明 Bott 定理。

\begin{proof}[Bott 定理的证明]
设 $P:\Gamma E\to\Gamma F$ 是椭圆齐性微分算子。因 $P$ 是 $G$-映射，它在每个有限维 $G$-不变子空间 $\Gamma_\alpha$ 上诱导线性映射 $P_\alpha:\Gamma_\alpha E\to\Gamma_\alpha F$。因 $P$ 椭圆，$\Ker P$ 与 $\Coker P$ 是有限维 $G$-模，Peter-Weyl 定理给出
\[
\Ker P=\bigoplus_\alpha\Ker P_\alpha,\qquad
\Coker P=\bigoplus_\alpha\Coker P_\alpha,
\]
其中除有限多个 $\alpha$ 外 $\Ker P_\alpha$ 与 $\Coker P_\alpha$ 为零。于是
\[
\Ind_G P=[\Ker P]-[\Coker P]=\sum_\alpha[\Ker P_\alpha]-[\Coker P_\alpha].
\]
又注意有有限维 $G$-模的正合列
\[
0\longrightarrow\Ker P_\alpha\longrightarrow\Gamma_\alpha E
\xrightarrow{\ P_\alpha\ }\Gamma_\alpha F\longrightarrow\Coker P_\alpha\longrightarrow 0,
\]
可对 $R(G)$ 使用由复形给出的 $K_G$ 理论构造。此复形的支撑为空，其欧拉示性数（即各 $G$-模的交错和）在 $R(G)$ 中必为零。于是
\[
\begin{aligned}
[\Ker P_\alpha]-[\Coker P_\alpha]
&=[\Gamma_\alpha E]-[\Gamma_\alpha F]\\
&=\ip{w_\alpha}{[\Gamma E]}_G w_\alpha=\ip{w_\alpha}{[\Gamma F]}_G w_\alpha.
\end{aligned}
\]
对 $\alpha$ 求和，得到
\[
\Ind_G P=\sum_\alpha\ip{w_\alpha}{[\Gamma E]}_G w_\alpha-\sum_\alpha\ip{w_\alpha}{[\Gamma F]}_G w_\alpha=[\Gamma E]-[\Gamma F],
\]
这正是所要的结论。
\end{proof}

\transnote{「支撑为空」是说该复形的象征在无穷远处可逆，故复形在 $K_G$ 中代表零类；等价地，它在 $R(G)$ 中的欧拉示性数 $\sum_i(-1)^i[\Gamma_\alpha^i]$ 为零。正是这一点把 $[\Ker P_\alpha]-[\Coker P_\alpha]$ 换成 $[\Gamma_\alpha E]-[\Gamma_\alpha F]$。}
